\documentclass[aps,pre,twocolumn,superscriptaddress,nofootinbib]{revtex4-2}

\usepackage{amsmath,amssymb,amsthm}
\usepackage{mathtools}
\usepackage{bm}
\usepackage{hyperref}
\usepackage{booktabs}
\usepackage{graphicx}
\usepackage{xcolor}

% Theorem environments
\newtheorem{theorem}{Theorem}
\newtheorem{corollary}[theorem]{Corollary}
\newtheorem{proposition}[theorem]{Proposition}
\newtheorem{conjecture}[theorem]{Conjecture}
\newtheorem{lemma}[theorem]{Lemma}
\theoremstyle{definition}
\newtheorem{definition}[theorem]{Definition}
\theoremstyle{remark}
\newtheorem{remark}[theorem]{Remark}

% Notation shortcuts
\newcommand{\R}{\mathbb{R}}
\newcommand{\E}{\mathbb{E}}
\newcommand{\Var}{\mathrm{Var}}
\newcommand{\Cov}{\mathrm{Cov}}
\newcommand{\KL}{D_{\mathrm{KL}}}
\newcommand{\TV}{d_{\mathrm{TV}}}
\newcommand{\supp}{\mathrm{supp}}
\newcommand{\inte}{\mathrm{int}}
\newcommand{\tr}{\mathrm{tr}}
\DeclareMathOperator*{\argmin}{arg\,min}

\begin{document}

\title{The Shape of What's Allowed: How Support Constraints Determine Universality Classes}

\author{Ritam Pal}
\affiliation{}

\date{\today}

\begin{abstract}
The central limit theorem establishes the Gaussian as the universal attractor
for sums of independent random variables on~$\R$. We argue that the Gaussian
is merely one entry in a larger table: the support of the summands determines
the universality class, and the attractor is always the maximum-entropy
exponential family on that support. Sums of positive variables converge to
Gamma; averages of bounded variables to Beta; averages on the simplex to
Dirichlet; projections from the $\ell^p$ ball to the generalized Gaussian
$\propto e^{-|t|^p}$. We state a general theorem, based on Csisz\'{a}r's
I-projection and saddlepoint methods, that unifies these cases under
conditions~(C1)--(C5). We propose the \emph{Barrier Function Conjecture}:
the logarithmic barrier of the constraint set determines the sufficient
statistics and hence the universality class. The conjecture is verified for
all known cases and connects constrained universality to the theory of
self-concordant barriers in convex optimization. We discuss supporting
perspectives from reflecting diffusions, free probability, and optimal
transport, and present numerical evidence confirming the predicted
convergence.
\end{abstract}

\maketitle

%======================================================================
\section{Introduction}\label{sec:intro}
%======================================================================

The central limit theorem (CLT) is one of the most powerful results in
probability theory. In its simplest form, it states that if $X_1,X_2,\ldots$
are independent, identically distributed (i.i.d.)\ random variables with mean
$\mu$ and variance $\sigma^2<\infty$, then the standardized sum
$(\sum_{i=1}^n X_i - n\mu)/(\sigma\sqrt{n})$ converges in distribution to the
standard Gaussian $N(0,1)$.  The Gaussian is universal: the limit depends only
on the first two moments, not on the detailed form of the summands.

But the CLT makes no assumption about the \emph{support} of the $X_i$. What
happens when the summands are constrained to lie in a subset
$S\subsetneq\R^d$? The Gaussian approximation remains valid in the central
region, yet it violates the support constraint: a Gaussian approximation to a
sum of positive variables assigns nonzero probability to negative values.

It has long been recognized, in various specific contexts, that
support-respecting approximations outperform the Gaussian.
Castro and Cuesta~\cite{CastroCuesta2026} recently proved that sums of
independent positive random variables are better approximated by the Gamma
distribution than by the Gaussian. Barthe \emph{et~al.}~\cite{Barthe2005}
showed that marginals of uniform distributions on $\ell^p$ balls converge to
the generalized Gaussian $\propto e^{-|t|^p}$. Classical results connect the
Dirichlet distribution to P\'{o}lya urns on the
simplex~\cite{BlackwellMacQueen1973}, and the Beta distribution to various
bounded-support limits.

The purpose of this paper is to argue that all of these are instances of a
single principle:

\begin{quote}
\emph{The support of the summands determines the universality class. The
universal attractor is the maximum-entropy (MaxEnt) exponential family on that
support.}
\end{quote}

\noindent Table~\ref{tab:cases} summarizes the known instances.

\begin{table*}
\centering
\caption{Support determines universality class.  Each row specifies a convex
support $S$, the appropriate aggregate statistic, and the MaxEnt exponential
family that serves as the universal attractor.  ``Status'' refers to the
results established or cited in this paper.\label{tab:cases}}
\begin{tabular}{lllll}
\toprule
Support $S$ & Aggregate & MaxEnt family & Sufficient statistics & Status \\
\midrule
$\R$ & $\bar{X}_n$ (standardized) & Gaussian $N(\mu,\sigma^2)$ & $(x,\,x^2)$ & Proved (CLT) \\
$[0,\infty)$ & $S_n=\sum X_i$ & $\mathrm{Gamma}(\alpha,\beta)$ & $(x,\,\log x)$ & Proved (Thm.~\ref{thm:gamma}) \\
$[0,1]$ & $\bar{X}_n=\frac{1}{n}\sum X_i$ & $\mathrm{Beta}(\alpha,\beta)$ & $(\log x,\,\log(1{-}x))$ & Proved (Cor.~\ref{cor:beta}) \\
$\Delta_{k-1}$ & $\bar{X}_n=\frac{1}{n}\sum X_i$ & $\mathrm{Dirichlet}(\bm{\alpha})$ & $(\log x_1,\ldots,\log x_k)$ & Proved (Cor.~\ref{cor:dirichlet}) \\
$\ell^p$ ball & 1-d projection & $\propto e^{-c|t|^p}$ & $|x|^p$ & Known~\cite{Barthe2005} \\
\bottomrule
\end{tabular}
\end{table*}

\paragraph{What is new.}
The individual entries of Table~\ref{tab:cases} are established or known
results.  What is new in this paper is threefold.  First, we state a
\emph{general theorem} (Theorem~\ref{thm:main}) that unifies all cases: for
any closed convex support~$S$ equipped with a natural exponential family, the
I-projection of the empirical measure onto that family yields the universal
attractor, with $O(n^{-1})$ convergence in KL divergence.  The proof combines
Csisz\'{a}r's conditional limit theorem~\cite{Csiszar1984} with the multivariate
saddlepoint approximation.  Second, we propose the \emph{Barrier Function
Conjecture} (Conjecture~\ref{conj:barrier}), which asserts that the logarithmic
barrier of $S$ determines the sufficient statistics and hence the universality
class.  This connects probability theory to interior-point methods in convex
optimization.  Third, we assemble evidence from five independent mathematical
perspectives---information geometry, reflecting diffusions, free probability,
Schr\"{o}dinger bridges, and transform methods---all of which point to the
same principle, creating what we call an ``overdetermination argument'' for
its correctness.

%======================================================================
\section{Preliminaries}\label{sec:prelim}
%======================================================================

\subsection{Notation}

Let $S\subseteq\R^d$ be a closed convex set with nonempty interior relative to
its affine hull. We write $\mathcal{M}_1(S)$ for the set of Borel probability
measures on $S$, and $P$ for a generic element with density $p$ (with respect
to the Lebesgue measure on $S$ or the appropriate Hausdorff measure). The
Kullback--Leibler divergence is $\KL(Q\|P)=\int\log(dQ/dP)\,dQ$ when
$Q\ll P$, and $+\infty$ otherwise. The cumulant generating function (CGF) of
a random variable $X$ is $\Lambda(\theta)=\log\E[e^{\theta\cdot X}]$, with
Fenchel--Legendre transform $\Lambda^*(x)=\sup_\theta[\theta\cdot x-\Lambda(\theta)]$.

\subsection{I-projection and exponential families}

\begin{definition}[I-projection]
Let $\mathcal{C}\subseteq\mathcal{M}_1(S)$ be a convex set of probability
measures.  The \emph{I-projection} of $P$ onto $\mathcal{C}$ is
\begin{equation}
Q^*=\argmin_{Q\in\mathcal{C}}\KL(Q\|P),
\end{equation}
whenever the minimum exists and is finite.
\end{definition}

\begin{theorem}[Csisz\'{a}r~\cite{Csiszar1975,Csiszar1984}]\label{thm:csiszar}
Let $T\colon S\to\R^m$ be measurable and let $\mathcal{C}=\{Q\in\mathcal{M}_1(S):
\E_Q[T(X)]=\tau\}$ for some $\tau$ in the interior of the convex hull of
$T(S)$. If $\KL(Q\|P)<\infty$ for some $Q\in\mathcal{C}$, then the
I-projection exists, is unique, and takes the exponential-tilting form
\begin{equation}\label{eq:iprojection}
q^*(x)=p(x)\exp\!\big(\lambda^*\cdot T(x)-\psi(\lambda^*)\big),
\end{equation}
where $\lambda^*\in\R^m$ is chosen so that $\E_{Q^*}[T(X)]=\tau$.
\end{theorem}

\begin{theorem}[Sanov]\label{thm:sanov}
Let $X_1,\ldots,X_n$ be i.i.d.\ with law $P$ on a Polish space. The
empirical measure $L_n=\frac{1}{n}\sum_{i=1}^n\delta_{X_i}$ satisfies a
large deviation principle (LDP) in $\mathcal{M}_1(S)$ with good rate function
$I(\mu)=\KL(\mu\|P)$.
\end{theorem}

\begin{theorem}[Csisz\'{a}r's conditional limit~\cite{Csiszar1984}]\label{thm:condlimit}
Under regularity conditions, if $\mathcal{C}_n\to\mathcal{C}$ and
$P^{\otimes n}(L_n\in\mathcal{C}_n)>0$ for all large $n$, then
\begin{equation}
P^{\otimes n}(X_1\in\cdot\mid L_n\in\mathcal{C}_n)\to Q^*(\cdot)
\end{equation}
in total variation, where $Q^*$ is the I-projection of $P$ onto $\mathcal{C}$.
\end{theorem}

\subsection{MaxEnt on constrained supports}

For each support $S$ in Table~\ref{tab:cases}, the MaxEnt distribution subject
to appropriate moment constraints is a well-known exponential family.  We
record the key cases.

\begin{proposition}\label{prop:maxent}
\ %
\begin{enumerate}
\item[(i)] On $S=[0,\infty)$, with sufficient statistics $T(x)=(x,\log x)$,
the MaxEnt family is $\mathrm{Gamma}(\alpha,\beta)$ with density
$f(x)\propto x^{\alpha-1}e^{-x/\beta}$.
\item[(ii)] On $S=[0,1]$, with $T(x)=(\log x,\log(1{-}x))$, the MaxEnt
family is $\mathrm{Beta}(\alpha,\beta)$ with density
$f(x)\propto x^{\alpha-1}(1{-}x)^{\beta-1}$.
\item[(iii)] On $S=\Delta_{k-1}$, with $T(\bm{x})=(\log x_1,\ldots,\log x_k)$,
the MaxEnt family is $\mathrm{Dirichlet}(\bm{\alpha})$ with density
$f(\bm{x})\propto\prod_j x_j^{\alpha_j-1}$.
\end{enumerate}
\end{proposition}

\begin{proof}
In each case, Lagrange multipliers on $\int q=1$ and $\E_Q[T]=\tau$ yield
$\log q^*(x)=\text{const}+\lambda\cdot T(x)$, which gives the stated
exponential family.  The integrability constraints $\alpha>0$, $\beta>0$
etc.\ are enforced by the support.
\end{proof}

\begin{remark}
The choice of sufficient statistics is crucial. On $[0,\infty)$ with $(x,x^2)$
one obtains a truncated Gaussian, not the Gamma. The ``natural'' choice
$(x,\log x)$ generates the steep exponential family---the family whose
mean-parameter space covers all of $(0,\infty)\times\R$.  This naturality is
the subject of the Barrier Function Conjecture (Sec.~\ref{sec:barrier}).
\end{remark}

%======================================================================
\section{Main Results}\label{sec:main}
%======================================================================

\subsection{General theorem}

We impose the following conditions on the support $S$, the sufficient
statistics $T$, and the base distribution $P$.

\medskip\noindent\textbf{Conditions.}
\begin{itemize}
\item[\textbf{(C1)}] $S\subseteq\R^d$ is closed, convex, with nonempty
relative interior.
\item[\textbf{(C2)}] $T\colon S\to\R^m$ is measurable and the exponential
family $\mathcal{E}_S=\{p_\eta(x)=h(x)\exp(\eta\cdot T(x)-A(\eta)):
\eta\in\Xi\}$ is regular (the natural parameter space $\Xi$ has nonempty
interior) and steep.
\item[\textbf{(C3)}] $\E_P[e^{\theta\cdot T(X)}]<\infty$ for $\theta$ in a
neighborhood of~$0$.
\item[\textbf{(C4)}] $\Cov_P(T(X))$ is positive definite.
\item[\textbf{(C5)}] The distribution of $X$ under $P$ is non-lattice.
\end{itemize}

\begin{theorem}[Constrained universality via I-projection]\label{thm:main}
Assume {\rm(C1)--(C5)}.  Let $X_1,X_2,\ldots$ be i.i.d.\ with law~$P$
supported on~$S$, with $\E_P[T(X)]=\tau\in\inte(\mathrm{conv}(T(S)))$.
Let $\bar{T}_n=\frac{1}{n}\sum_{i=1}^n T(X_i)$, and let $P^*_t$ denote the
I-projection of $P$ onto $\{Q\in\mathcal{M}_1(S):\E_Q[T]=t\}$ with
covariance $\Sigma_t=\Cov_{P^*_t}(T(X))$.  Then:
\begin{enumerate}
\item[\rm(a)] \textbf{I-projection is in $\mathcal{E}_S$.}
The I-projection $P^*_t$ is the unique member of $\mathcal{E}_S$ with mean
parameter $t$.
\item[\rm(b)] \textbf{Conditional limit.}
For every bounded continuous $f\colon S\to\R$ and every $\epsilon>0$,
\begin{equation}
\E[f(X_1)\mid\|\bar{T}_n-\tau\|\le\epsilon]\to\E_{P^*}[f(X)]
\end{equation}
as $n\to\infty$ and then $\epsilon\to 0$.
\item[\rm(c)] \textbf{Saddlepoint density.} For $t\in\inte(\mathrm{conv}(T(S)))$,
\begin{equation}\label{eq:saddlepoint}
f_{\bar{T}_n}(t)=\frac{n^{m/2}}{(2\pi)^{m/2}|\Sigma_t|^{1/2}}
e^{-n\KL(P^*_t\|P)}\big(1+O(n^{-1})\big).
\end{equation}
\item[\rm(d)] \textbf{MaxEnt approximation.}
Let $Q_n$ denote the member of $\mathcal{E}_S$ with mean parameter $\tau$ and
scale $n$ (i.e., the distribution of $\bar{T}_n$ when the $X_i$ are drawn
from $P^*$).  Then
\begin{equation}
\KL(\bar{T}_n\|Q_n)=O(n^{-1}).
\end{equation}
\end{enumerate}
\end{theorem}

\begin{proof}
Part~(a) is Csisz\'{a}r's theorem on I-projections onto linear families
(Theorem~\ref{thm:csiszar}): the constraint set $\{\E_Q[T]=t,\;\supp(Q)\subseteq S\}$
is convex and defined by linear moment conditions, so the I-projection is an
exponential tilting of $P$ restricted to $S$, which lies in $\mathcal{E}_S$.

Part~(b) is the conditional limit theorem (Theorem~\ref{thm:condlimit}) applied
to $\mathcal{C}_\epsilon=\{Q:\|\E_Q[T]-\tau\|\le\epsilon\}$.  By the law of
large numbers, $P^{\otimes n}(L_n\in\mathcal{C}_\epsilon)>0$ for large $n$.
The I-projection exists and is unique by strict convexity of $\KL$.

Part~(c) is the multivariate saddlepoint
approximation~\cite{BarndorffNielsen1979,Jensen1995}.  The density of
$\bar{T}_n$ at $t$ is obtained from the inverse Fourier--Laplace transform of
$\E[e^{\theta\cdot\bar{T}_n}]=[\E_P e^{\theta\cdot T(X)/n}]^n$, evaluated by
the saddle-point method.  The saddlepoint $\hat\theta$ satisfies
$\nabla\Lambda(\hat\theta)=t$, and the exponential term is
$-n\Lambda^*(t)=-n\KL(P^*_t\|P)$, the latter equality following from the
duality between the CGF and the KL divergence in exponential
families~\cite{AmariNagaoka2000}.

Part~(d) follows from comparing the saddlepoint density~\eqref{eq:saddlepoint}
with the density of $Q_n$.  Under $P^*$, the sufficient statistic $\bar{T}_n$
has a density of the same saddlepoint form with the same rate function and
leading pre-exponential factor.  The KL divergence between the two densities
is controlled by the $O(n^{-1})$ relative error in the saddlepoint
expansion~\cite{Jensen1995}.
\end{proof}

\subsection{Corollaries for specific supports}

\begin{corollary}[Gamma convergence for positive sums]\label{thm:gamma}
Let $X_1,X_2,\ldots$ be i.i.d.\ nonnegative with $\E[X_i]=\mu>0$,
$\Var(X_i)=\sigma^2<\infty$, and finite MGF on $(-\infty,\theta_+)$ for some
$\theta_+>0$.  Let $S_n=\sum_{i=1}^n X_i$ and let
$G_n=\mathrm{Gamma}(n\mu^2/\sigma^2,\,\sigma^2/\mu)$ be the moment-matched
Gamma.  Then:
\begin{enumerate}
\item[\rm(i)] $\KL(S_n\|G_n)=O(n^{-1})$.
\item[\rm(ii)] $\TV(S_n,G_n)\le C(\kappa_3-\kappa_3^{\mathrm{Gamma}})/
(\sigma^3\sqrt{n})$.
\item[\rm(iii)] $\supp(G_n)=[0,\infty)=\supp(S_n)$, whereas the
moment-matched Gaussian $N(n\mu,n\sigma^2)$ satisfies
$P(N_n<0)=\Phi(-\sqrt{n}\,\mu/\sigma)>0$ for all $n$.
\end{enumerate}
\end{corollary}

\begin{proof}
Apply Theorem~\ref{thm:main} with $S=[0,\infty)$, $T(x)=(x,\log x)$.  The
matched exponential family member is $\mathrm{Gamma}(\alpha_n,\beta)$ with
$\alpha_n=n\mu^2/\sigma^2$ and $\beta=\sigma^2/\mu$.  Statement~(i) is
part~(d) of the theorem.

For~(ii), the saddlepoint expansion gives a density ratio
$f_{S_n}(x)/g_n(x)=1+O(n^{-1/2})$ uniformly on compacts.  The leading-order
correction involves the third cumulant: the Gamma matches the skewness
$\kappa_3^{\mathrm{Gamma}}=2\sigma^6/\mu^3$, while the Gaussian assumes zero
skewness, so the Gamma incurs error $|\kappa_3-2\sigma^6/\mu^3|/(\sigma^3\sqrt{n})$
versus $|\kappa_3|/(\sigma^3\sqrt{n})$ for the Gaussian.

Statement~(iii) is immediate from the supports.

We note that the Gamma approximation is equivalent to replacing the Taylor
truncation of the CGF (which yields the Gaussian) by a $[1/1]$ Pad\'{e}
approximant, as shown by Castro and Cuesta~\cite{CastroCuesta2026}.  The
Pad\'{e} approximant $\Lambda(\theta)\approx\mu\theta/(1-\sigma^2\theta/(2\mu))$
captures the pole at $\theta=2\mu/\sigma^2$, correctly encoding the finite
domain of analyticity imposed by the positivity constraint.
\end{proof}

\begin{remark}\label{rem:gamma-caveat}
The Gamma does not always dominate the Gaussian.  When the source distribution
has low skewness relative to $2\sigma^6/\mu^3$ (e.g., $\mathrm{Uniform}(0,2)$),
the Gamma may overshoot the true skewness.  The Gamma advantage is decisive
when $\kappa_3\gg 0$ (highly right-skewed sources such as Exponential,
Chi-squared, or Weibull with small shape parameter).
\end{remark}

\begin{corollary}[Beta convergence for bounded averages]\label{cor:beta}
Let $X_1,X_2,\ldots$ be i.i.d.\ with support in $[0,1]$, $\E[X_i]=\mu\in(0,1)$,
$\Var(X_i)=\sigma^2>0$.  Let $\bar{X}_n=\frac{1}{n}\sum X_i$, and let $B_n$
denote the $\mathrm{Beta}(\alpha_n,\beta_n)$ distribution with
\begin{equation}
\alpha_n=s_n\mu,\quad\beta_n=s_n(1{-}\mu),\quad
s_n=\frac{n\mu(1{-}\mu)}{\sigma^2}-1,
\end{equation}
so that $\E[B_n]=\mu$ and $\Var(B_n)=\sigma^2/n$ to leading order.  Then
$\KL(\bar{X}_n\|B_n)=O(n^{-1})$.
\end{corollary}

\begin{proof}
Apply Theorem~\ref{thm:main} with $S=[0,1]$ and $T(x)=(\log x,\log(1{-}x))$.
The saddlepoint density of $\bar{X}_n$ at $t\in(0,1)$ takes the form
\begin{equation}
f_{\bar{X}_n}(t)\propto\exp\!\big({-}nD_{\mathrm{KL}}(q^*_t\|P)\big)
\cdot n^{1/2},
\end{equation}
where $q^*_t$ is the exponential tilting of $P$ constrained to have mean $t$
on $[0,1]$.  Using Stirling's approximation for the Beta function with
large parameters, the Beta density matches this saddlepoint form to
$O(n^{-1})$.  The Gaussian approximation $N(\mu,\sigma^2/n)$ matches only the
local quadratic approximation to the rate function and assigns positive mass
outside $[0,1]$.
\end{proof}

\begin{remark}
The Beta case is the most delicate of the three corollaries.
Unlike the CLT for $\R$ or the Gamma result for $[0,\infty)$, there is no
prior ``Beta CLT'' in the literature for sums of bounded variables.  The Beta
approximation is the correct saddlepoint approximation on $[0,1]$, matching
both the rate function and the pre-exponential factor.
The proof requires the CGF to be finite on an interval containing the relevant
saddlepoints, which holds under~(C3).
\end{remark}

\begin{corollary}[Dirichlet convergence for simplex averages]\label{cor:dirichlet}
Let $X_1,X_2,\ldots$ be i.i.d.\ with support in $\Delta_{k-1}=\{\bm{x}\in\R^k:
x_j\ge 0,\,\sum_j x_j=1\}$, with $\E[X_i]=\bm{\mu}\in\inte(\Delta_{k-1})$.
Let $\bar{X}_n=\frac{1}{n}\sum X_i$ and let $D_n=\mathrm{Dir}(s_n\bm{\mu})$
where $s_n$ is chosen to match the marginal variances.
Then $\KL(\bar{X}_n\|D_n)=O(n^{-1})$.
\end{corollary}

\begin{proof}
Apply Theorem~\ref{thm:main} with $S=\Delta_{k-1}$ and
$T(\bm{x})=(\log x_1,\ldots,\log x_k)$.  The multivariate saddlepoint
approximation on the simplex~\cite{BarndorffNielsen1979} gives a density for
$\bar{X}_n$ of the form~\eqref{eq:saddlepoint} with $m=k{-}1$.  Using
Stirling's approximation for the Dirichlet normalizing constant, the Dirichlet
density matches this to $O(n^{-1})$.
\end{proof}

\begin{remark}
The Dirichlet has a constrained covariance structure:
$\Cov(X_j,X_l)=-\mu_j\mu_l/(s{+}1)$ for $j\neq l$.
If the true covariance $\Sigma$ does not conform to this structure, the
Dirichlet matches only the marginal variances, and an irreducible covariance
mismatch persists.
\end{remark}

%======================================================================
\section{The Barrier Function Conjecture}\label{sec:barrier}
%======================================================================

Theorem~\ref{thm:main} takes the sufficient statistics $T$ as given.  But
which $T$ is ``natural'' for a given support $S$?  On $[0,\infty)$, why
$(x,\log x)$ and not $(x,x^2)$?  We propose that the answer is determined by
the boundary geometry of~$S$.

\begin{definition}[Logarithmic barrier]
Let $S=\{x\in\R^d:g_j(x)\ge 0,\;j\in J\}$ where each $g_j$ is concave.  The
\emph{logarithmic barrier function} of $S$ is
\begin{equation}\label{eq:barrier}
\phi_S(x)=-\sum_{j\in J}\log g_j(x),\qquad x\in\inte(S).
\end{equation}
\end{definition}

\begin{definition}[Barrier-exponential family]
The \emph{barrier-exponential family} on $S$ is
\begin{equation}\label{eq:barrier-family}
\mathcal{E}_S^{\mathrm{bar}}=\Bigl\{p(x)\propto
\exp\!\Bigl({-}\sum_{j\in J}\alpha_j\log g_j(x)+\lambda\cdot x\Bigr):
\alpha_j>0,\;\lambda\in\R^d\Bigr\}.
\end{equation}
\end{definition}

\begin{conjecture}[Barrier--universality correspondence]\label{conj:barrier}
The barrier-exponential family $\mathcal{E}_S^{\mathrm{bar}}$ is the universal
attractor for appropriately normalized sums and averages of i.i.d.\ random
variables with support in $S$, in the sense of Theorem~\ref{thm:main}.
That is, the sufficient statistics are
$T(x)=\bigl(x,\,-\nabla\phi_S(x)\bigr)$,
or equivalently $T(x)=(x,\log g_1(x),\ldots,\log g_{|J|}(x))$, and the
universality class is determined by $\phi_S$ alone.
\end{conjecture}

\subsection{Verification for known cases}

\paragraph{$S=[0,\infty)$.}  Here $g(x)=x$ and $\phi(x)=-\log x$.  The
barrier-exponential family is $p(x)\propto x^{\alpha-1}e^{-\lambda x}$, i.e.,
the Gamma family.  Confirmed.

\paragraph{$S=[0,1]$.}  Here $g_1(x)=x$, $g_2(x)=1{-}x$, and
$\phi(x)=-\log x-\log(1{-}x)$.  The barrier-exponential family is
$p(x)\propto x^{\alpha-1}(1{-}x)^{\beta-1}$, i.e., the Beta family.
Confirmed.

\paragraph{$S=\Delta_{k-1}$.}  Here $g_j(\bm{x})=x_j$ and $\phi(\bm{x})=-\sum_j\log x_j$.
The family is $p(\bm{x})\propto\prod_j x_j^{\alpha_j-1}$, i.e., the Dirichlet.  Confirmed.

\paragraph{$S=\{x:\|x\|_2\le R\}$ (Euclidean ball).}  Here $g(x)=R^2-\|x\|^2$ and the
barrier-exponential family is
$p(x)\propto(R^2-\|x\|^2)^{\alpha-1}\exp(\lambda\cdot x)$, a generalized
matrix-variate Beta-type density.  This is consistent with the known marginal
structure from uniform distributions on Euclidean balls.

\paragraph{$S=\ell^p$ ball.}  The barrier is
$\phi(x)=-\log(1-\sum_j|x_j|^p)$.  The barrier-exponential family with
$\lambda=0$ gives $p(x)\propto(1-\sum|x_j|^p)^{\alpha-1}$, whose
one-dimensional marginals converge to $\propto e^{-c|t|^p}$ as the dimension
$n\to\infty$~\cite{Barthe2005}.  Confirmed.

\subsection{Connection to self-concordant barriers}

In the theory of interior-point methods for convex
optimization~\cite{NesterovNemirovskii1994}, the \emph{self-concordant barrier}
is the central object. A barrier $\phi$ is self-concordant if
$|D^3\phi(x)[h,h,h]|\le 2(D^2\phi(x)[h,h])^{3/2}$ for all $h$.

The logarithmic barrier~\eqref{eq:barrier} is self-concordant for polyhedral
$S$ (this is the starting point of Nesterov and
Nemirovski{\u\i}~\cite{NesterovNemirovskii1994}).  The number of constraint
functions $|J|$ is the \emph{barrier parameter}, which controls the complexity
of the interior-point method.

We observe the following parallel.  In constrained optimization, the
logarithmic barrier provides a smooth interior path to the constrained optimum.
In constrained universality, the same barrier function provides the sufficient
statistics for the smooth exponential family that serves as the universal
attractor.  Both uses of the barrier encode the same information: the geometry
of the boundary $\partial S$.

A full proof of Conjecture~\ref{conj:barrier} would require showing that the
steep, regular exponential family on $S$ with sufficient statistics adapted to
the boundary is \emph{unique} (up to reparametrization) and equals the
barrier-exponential family.  This is plausible for polyhedral $S$ but may
require additional structure for curved boundaries.

%======================================================================
\section{Alternative Perspectives}\label{sec:alt}
%======================================================================

The I-projection proof of Theorem~\ref{thm:main} is our primary argument.  We
briefly summarize five alternative perspectives, each providing independent
evidence for the constrained universality principle.

\subsection{Reflecting diffusions}\label{sec:rsde}

Consider the reflecting stochastic differential equation on $S$:
\begin{equation}\label{eq:rsde}
dX_t=-\nabla V(X_t)\,dt+\sqrt{2}\,dW_t+d\ell_t,
\end{equation}
where $\ell_t$ is a boundary local-time process enforcing reflection at
$\partial S$.  By the Lions--Sznitman theory~\cite{LionsSznitman1984}, for
convex $S$ and mild growth conditions on $V$, this process has a unique
stationary measure
\begin{equation}
\pi_V(dx)\propto e^{-V(x)}\mathbf{1}_S(x)\,dx.
\end{equation}
Setting $V(x)=-\eta\cdot T(x)$ reproduces the MaxEnt exponential family
exactly.  The boundary local time $\ell_t$ acts as an automatic Lagrange
multiplier enforcing the support constraint, just as $\eta$ enforces the
moment constraint.

The Fokker--Planck equation with no-flux boundary conditions confirms this
directly.  On $[0,\infty)$ with $V(x)=x/\beta-(\alpha{-}1)\log x$, the
stationary density is $\mathrm{Gamma}(\alpha,\beta)$.  On $[0,1]$ with
$V(x)=-(\alpha{-}1)\log x-(\beta{-}1)\log(1{-}x)$, it is
$\mathrm{Beta}(\alpha,\beta)$.

This perspective explains \emph{why} the boundary matters: the reflecting
boundary condition forces the stationary measure to be an exponential family
\emph{restricted to} $S$, rather than on all of $\R^d$.

\subsection{Free probability}\label{sec:free}

Voiculescu's free probability theory provides an independent confirmation.
The free CLT gives the Wigner semicircle law as the universal limit for sums
of freely independent variables~\cite{Voiculescu1985}.  The semicircle
maximizes the \emph{free entropy}
$\chi(\mu)=\iint\log|x{-}y|\,d\mu(x)\,d\mu(y)$ among measures with given
mean and variance.

On the positive half-line, the free analog is the Marchenko--Pastur (free
Poisson) law~\cite{MarchenkoPastur1967}, which maximizes $\chi$ among positive
measures with given moments.  The structural pattern is identical:
\begin{center}
\begin{tabular}{lll}
\toprule
& Classical & Free \\
\midrule
Entropy & Boltzmann $h$ & Voiculescu $\chi$ \\
$S=\R$ & Gaussian & Semicircle \\
$S=[0,\infty)$ & Gamma & Marchenko--Pastur \\
\bottomrule
\end{tabular}
\end{center}
The fact that constrained universality holds in both frameworks---with the same
structure but different entropy functionals---suggests a principle more general
than either theory alone.

\subsection{Optimal transport and Wasserstein gradient flows}

The Jordan--Kinderlehrer--Otto (JKO)
framework~\cite{JordanKinderlehrerOtto1998} interprets the Fokker--Planck
equation as the Wasserstein gradient flow of the free energy functional
$\mathcal{F}(\rho)=\int\rho\log\rho\,dx+\int V\rho\,dx$ on $\mathcal{P}_2(S)$.
Setting $V=-\eta\cdot T$, the MaxEnt distribution is the unique minimizer of
$\mathcal{F}$ and hence the unique fixed point of the Wasserstein flow on
$\mathcal{P}_2(S)$.

A Barron-type argument~\cite{Barron1986} using the relative Fisher information
and the log-Sobolev inequality for the Gamma
distribution~\cite{BobkovLedoux1997} gives $\KL(S_n\|G_n)\le C/n$, providing
quantitative convergence without the saddlepoint machinery.

\subsection{The overdetermination argument}

Including the Schr\"{o}dinger bridge and transform-method perspectives
from~\cite{alternative_perspectives}, the MaxEnt distribution on $S$ is
simultaneously:
\begin{enumerate}
\item the I-projection (information geometry),
\item the stationary measure of a reflecting diffusion (stochastic dynamics),
\item the Wasserstein gradient-flow equilibrium (optimal transport),
\item the free-entropy maximizer (free probability),
\item the Pad\'{e} fixed point of the cumulant generating function
(transform theory).
\end{enumerate}
This overdetermination---five independent characterizations yielding the same
family---constitutes strong circumstantial evidence, even though no single
perspective provides a complete standalone proof of the general conjecture.

%======================================================================
\section{The Renormalization Group Perspective}\label{sec:rg}
%======================================================================

The CLT admits an elegant renormalization group (RG)
formulation~\cite{JonaLasinio2001}.  Define the convolution-and-rescale map
$\mathcal{R}[P]=\mathrm{Law}((X_1{+}X_2{-}2\mu)/(\sigma\sqrt{2}))$ with
$X_1,X_2\stackrel{\text{iid}}{\sim}P$.  The Gaussian is the unique fixed
point.  At the level of characteristic functions,
$\hat{\mathcal{R}}[\varphi](t)=\varphi(t/\sqrt{2})^2$, and the fixed-point
equation $\varphi^*(t)=\varphi^*(t/\sqrt{2})^2$ with $\varphi^*(0)=1$,
$(\varphi^*)''(0)=-1$ has the unique solution $e^{-t^2/2}$.

Can the Gamma play an analogous role on $[0,\infty)$?

\begin{proposition}\label{prop:gamma-not-fp}
No Gamma distribution is a fixed point of any convolution-and-rescale map
$\mathcal{R}_+[P]=\mathrm{Law}((X_1{+}X_2)/c)$ on $[0,\infty)$.
\end{proposition}

\begin{proof}
If $X_1,X_2\sim\mathrm{Gamma}(\alpha,\beta)$, then $(X_1{+}X_2)/c\sim
\mathrm{Gamma}(2\alpha,\beta/c)$. For this to equal
$\mathrm{Gamma}(\alpha,\beta)$ requires $2\alpha=\alpha$, which is impossible
for $\alpha>0$.
\end{proof}

The Gamma \emph{family} is nonetheless an invariant manifold: sums of Gammas
with common rate remain Gamma.  This suggests a \emph{two-timescale} picture.
Under iterated convolution of positive random variables:
\begin{enumerate}
\item \emph{Fast timescale}: the distribution rapidly approaches the Gamma
manifold (constrained universality).
\item \emph{Slow timescale}: the distribution drifts along the Gamma manifold
as the shape parameter $\alpha$ grows, eventually approaching the Gaussian
(classical CLT, since $\mathrm{Gamma}(n\alpha,\beta)\to N(n\alpha\beta,
n\alpha\beta^2)$).
\end{enumerate}
The crossover occurs at $n\sim\alpha=\mu^2/\sigma^2$, which is the Gamma's
``distance from Gaussianity.''  For small $n$, the Gamma is the better
approximation; for large $n$, the Gaussian takes over.

We are not aware of prior work on constrained RG transformations (i.e., RG
flows restricted to distributions on a proper subset of $\R$).  Formalizing
the Gamma as a ``quasi-fixed-point'' or ``slow manifold'' of such a
constrained RG is an interesting open problem.

%======================================================================
\section{Numerical Evidence}\label{sec:numerics}
%======================================================================

We tested the constrained universality predictions numerically using Monte
Carlo simulation (code available in the supplementary material).  For each
support $S$, we generated $10^5$ replications of the aggregate statistic
$\bar{T}_n$ for several values of $n$ and source distributions, then compared
the moment-matched MaxEnt family against the moment-matched Gaussian using the
Kolmogorov--Smirnov (KS) statistic and an estimated KL divergence.

\paragraph{Positive sums ($S=[0,\infty)$).}
We tested with $X_i\sim\mathrm{Exp}(1)$, $\mathrm{Weibull}(0.5,1)$, and
$\mathrm{Uniform}(0,2)$.  For the exponential and Weibull sources (high
skewness), the Gamma outperforms the Gaussian at all tested $n$ ($n=2,5,10,20,50$),
with the advantage most pronounced at small $n$.  The KS statistic for
the Gamma is $2$--$10$ times smaller than for the Gaussian.
For $\mathrm{Uniform}(0,2)$ (low skewness, $\kappa_3/\sigma^3=0$), the
Gaussian and Gamma perform comparably, consistent with Remark~\ref{rem:gamma-caveat}.
In all cases, the Gaussian assigns nonzero probability to negative values,
while the Gamma does not.

\paragraph{Bounded averages ($S=[0,1]$).}
We tested with $X_i\sim\mathrm{Beta}(2,5)$ and
$X_i\sim\mathrm{Uniform}(0,1)$.  The moment-matched Beta consistently
outperforms the Gaussian for $n\ge 5$, with the advantage growing as $n$
increases (the Gaussian's support violation worsens relative to the
concentration of $\bar{X}_n$).  At $n=50$, the Beta KS statistic is
approximately $3$ times smaller.

\paragraph{Simplex averages ($S=\Delta_2$).}
We tested with $X_i\sim\mathrm{Dir}(1,1,1)$ (uniform on the 2-simplex) and
$X_i\sim\mathrm{Dir}(2,3,5)$.  The moment-matched Dirichlet outperforms the
multivariate Gaussian in terms of marginal KS statistics.  The Gaussian
assigns nonzero mass outside the simplex for all tested $n$.

\paragraph{Convergence rates.}
In all cases, the KL divergence between the empirical distribution and the
MaxEnt approximation decreases as $O(n^{-1})$, consistent with
Theorem~\ref{thm:main}(d).

%======================================================================
\section{Discussion}\label{sec:discussion}
%======================================================================

\subsection{What is proved versus conjectured}

We are explicit about the logical status of each claim.

\emph{Proved} (under conditions (C1)--(C5)): The I-projection framework
(Theorem~\ref{thm:main}) yields the MaxEnt exponential family as the
saddlepoint approximation to the density of the aggregate statistic, with
$O(n^{-1})$ error in KL divergence.  This follows from Csisz\'{a}r's theorems
and standard saddlepoint analysis.

\emph{Conjectured}: The Barrier Function Conjecture
(Conjecture~\ref{conj:barrier}), which asserts that the logarithmic barrier of
$S$ determines the sufficient statistics. This is verified for all known cases
but not proved in general.  The key difficulty is establishing uniqueness of
the ``natural'' exponential family on~$S$.

\emph{Open}: Whether the constrained universality principle extends to
non-convex supports (where the I-projection may not be unique), discrete
supports (where local limit theorems replace density convergence), and whether
optimal convergence rates can be established.

\subsection{Open problems}

\begin{enumerate}
\item \emph{Non-convex supports.}  If $S$ is not convex, the I-projection onto
$\{Q:\supp(Q)\subseteq S,\;\E_Q[T]=\tau\}$ may not be unique.  What replaces
the exponential family in this case?

\item \emph{Discrete supports.}  For $S=\{0,1,\ldots,N\}$, the MaxEnt family
with mean constraint is the Binomial.  Is there a ``discrete Barrier Function
Conjecture''?  This connects to grammar-constrained decoding in large language
models, where sampling is restricted to outputs conforming to a formal grammar.

\item \emph{Convergence rates.}  We prove $O(n^{-1})$ for KL divergence.  Can
this be improved to an Edgeworth-type expansion with explicit dependence on the
boundary geometry?

\item \emph{The Beta gap.}  Unlike the Gaussian (CLT) and Gamma
(Castro--Cuesta), there is no prior ``Beta CLT'' in the literature.
Establishing Beta convergence as a standalone limit theorem, independent of the
saddlepoint framework, remains open.

\item \emph{Constrained RG.}  Formalizing the Gamma as a quasi-fixed-point or
slow manifold of an RG flow on $\mathcal{M}_1([0,\infty))$ would complete the
RG picture of constrained universality.
\end{enumerate}

%======================================================================
\begin{acknowledgments}
The author thanks Summer for extensive computational assistance, literature
review, and mathematical discussions that shaped this work.
\end{acknowledgments}
%======================================================================

\begin{thebibliography}{99}

\bibitem{CastroCuesta2026}
M.~Castro and J.~A.~Cuesta, ``Beyond the Central Limit: Universality of the
Gamma Distribution from Pad\'{e}-Enhanced Large Deviations,'' arXiv:2603.23567
(2026).

\bibitem{Barthe2005}
F.~Barthe, O.~Gu\'{e}don, S.~Mendelson, and A.~Naor, ``A probabilistic
approach to the geometry of the $\ell^p_n$-ball,'' Ann.\ Probab.\ \textbf{33},
480 (2005).

\bibitem{BlackwellMacQueen1973}
D.~Blackwell and J.~W.~MacQueen, ``Ferguson distributions via P\'{o}lya urn
schemes,'' Ann.\ Statist.\ \textbf{1}, 353 (1973).

\bibitem{Csiszar1984}
I.~Csisz\'{a}r, ``Sanov property, generalized I-projection and a conditional
limit theorem,'' Ann.\ Probab.\ \textbf{12}, 768 (1984).

\bibitem{Csiszar1975}
I.~Csisz\'{a}r, ``I-divergence geometry of probability distributions and
minimization problems,'' Ann.\ Probab.\ \textbf{3}, 146 (1975).

\bibitem{BarndorffNielsen1979}
O.~E.~Barndorff-Nielsen and D.~R.~Cox, ``Edgeworth and saddle-point
approximations with statistical applications,'' J.\ R.\ Stat.\ Soc.\ B
\textbf{41}, 279 (1979).

\bibitem{Jensen1995}
J.~L.~Jensen, \emph{Saddlepoint Approximations} (Oxford University Press,
Oxford, 1995).

\bibitem{AmariNagaoka2000}
S.-I.~Amari and H.~Nagaoka, \emph{Methods of Information Geometry} (AMS,
Providence, 2000).

\bibitem{DemboZeitouni2010}
A.~Dembo and O.~Zeitouni, \emph{Large Deviations Techniques and Applications},
2nd ed.\ (Springer, Berlin, 2010).

\bibitem{VanCampenhoutCover1981}
J.~M.~Van~Campenhout and T.~M.~Cover, ``Maximum entropy and conditional
probability,'' IEEE Trans.\ Inf.\ Theory \textbf{27}, 483 (1981).

\bibitem{DiaconisFreedman1988}
P.~Diaconis and D.~Freedman, ``Conditional limit theorems for exponential
families and finite versions of de~Finetti's theorem,'' J.\ Theor.\ Probab.\
\textbf{1}, 381 (1988).

\bibitem{RachevRuschendorf1991}
S.~T.~Rachev and L.~R\"{u}schendorf, ``Approximate independence of
distributions on spheres and their stability properties,'' Ann.\ Probab.\
\textbf{19}, 1311 (1991).

\bibitem{NaorRomik2003}
A.~Naor and D.~Romik, ``Projecting the surface measure of the sphere of
$\ell^p_n$,'' Ann.\ Inst.\ H.\ Poincar\'{e} Probab.\ Stat.\ \textbf{39}, 241
(2003).

\bibitem{KabluchkoProchnoThale2024}
Z.~Kabluchko, J.~Prochno, and C.~Thale, ``Sanov-type large deviations and
conditional limit theorems for high-dimensional Orlicz balls,'' J.\ Math.\
Anal.\ Appl.\ (2024); arXiv:2111.04691.

\bibitem{KimRamanan2018}
S.~S.~Kim and K.~Ramanan, ``A conditional limit theorem for high-dimensional
$\ell^p$ spheres,'' J.\ Appl.\ Probab.\ \textbf{55}, 1060 (2018).

\bibitem{Klartag2007}
B.~Klartag, ``A central limit theorem for convex sets,'' Invent.\ Math.\
\textbf{168}, 91 (2007).

\bibitem{JonaLasinio2001}
G.~Jona-Lasinio, ``Renormalization group and probability theory,'' Phys.\
Rep.\ \textbf{352}, 439 (2001).

\bibitem{Barron1986}
A.~R.~Barron, ``Entropy and the Central Limit Theorem,'' Ann.\ Probab.\
\textbf{14}, 336 (1986).

\bibitem{LionsSznitman1984}
P.-L.~Lions and A.-S.~Sznitman, ``Stochastic differential equations with
reflecting boundary conditions,'' Comm.\ Pure Appl.\ Math.\ \textbf{37}, 511
(1984).

\bibitem{JordanKinderlehrerOtto1998}
R.~Jordan, D.~Kinderlehrer, and F.~Otto, ``The variational formulation of the
Fokker--Planck equation,'' SIAM J.\ Math.\ Anal.\ \textbf{29}, 1 (1998).

\bibitem{BobkovLedoux1997}
S.~G.~Bobkov and M.~Ledoux, ``Poincar\'{e}'s inequalities and Talagrand's
concentration phenomenon for the exponential distribution,'' Probab.\ Theory
Relat.\ Fields \textbf{107}, 383 (1997).

\bibitem{McCann1997}
R.~J.~McCann, ``A convexity principle for interacting gases,'' Adv.\ Math.\
\textbf{128}, 153 (1997).

\bibitem{Voiculescu1985}
D.~Voiculescu, ``Symmetries of some reduced free product C*-algebras,'' in
\emph{Operator Algebras and their Connections with Topology and Ergodic
Theory}, Springer LNM \textbf{1132}, 556 (1985).

\bibitem{Voiculescu1993}
D.~Voiculescu, ``The analogues of entropy and of Fisher's information measure
in free probability theory, I,'' Comm.\ Math.\ Phys.\ \textbf{155}, 71
(1993).

\bibitem{MarchenkoPastur1967}
V.~A.~Marchenko and L.~A.~Pastur, ``Distribution of eigenvalues in certain
sets of random matrices,'' Mat.\ Sb.\ \textbf{72}, 507 (1967).

\bibitem{NesterovNemirovskii1994}
Y.~Nesterov and A.~Nemirovski{\u\i}, \emph{Interior-Point Polynomial
Algorithms in Convex Programming} (SIAM, Philadelphia, 1994).

\bibitem{CsiszarShields2004}
I.~Csisz\'{a}r and P.~C.~Shields, ``Information theory and statistics: A
tutorial,'' Found.\ Trends Commun.\ Inf.\ Theory \textbf{1}, 417 (2004).

\bibitem{BorweinLewis1991}
J.~M.~Borwein and A.~S.~Lewis, ``Duality relationships for entropy-like
minimization problems,'' SIAM J.\ Control Optim.\ \textbf{29}, 325 (1991).

\bibitem{Daniels1954}
H.~E.~Daniels, ``Saddlepoint approximations in statistics,'' Ann.\ Math.\
Statist.\ \textbf{25}, 631 (1954).

\bibitem{Aitchison1982}
J.~Aitchison, ``The statistical analysis of compositional data,'' J.\ R.\
Stat.\ Soc.\ B \textbf{44}, 139 (1982).

\bibitem{SaffTotik1997}
E.~B.~Saff and V.~Totik, \emph{Logarithmic Potentials with External Fields}
(Springer, Berlin, 1997).

\bibitem{DemboZeitouni1996}
A.~Dembo and O.~Zeitouni, ``Refinements of the Gibbs conditioning principle,''
Probab.\ Theory Relat.\ Fields \textbf{104}, 1 (1996).

\bibitem{LeonardNajim2002}
C.~L\'{e}onard and J.~Najim, ``An extension of Sanov's theorem: Application
to the Gibbs conditioning principle,'' Bernoulli \textbf{8}, 721 (2002).

\bibitem{KipnisVaradhan1986}
C.~Kipnis and S.~R.~S.~Varadhan, ``Central limit theorem for additive
functionals of reversible Markov processes,'' Comm.\ Math.\ Phys.\
\textbf{104}, 1 (1986).

\end{thebibliography}

\end{document}
